\section{Extensional Foundations: The Mathematical Substrate}

\begin{definition}[Extensional Relational Structure]
	An extensional relational structure (or a conceptualization according to Genesereth and Nilsson) is a tuple $(D, \mathbf{R})$ where:
	\begin{itemize}
		\item $D$ is a set called the universe of discourse
		\item $\mathbf{R}$ is a set of relations on $D$
	\end{itemize}
\end{definition}

\subsection{Philosophical Context and Systemic Interpretations}

This deceptively simple mathematical structure conceals profound ontological choices. The extensional relational structure provides what we might term the \emph{phenomenal cross-section} of a domain—a snapshot of entities and their actual interrelations at a particular state. Yet this very extensionality creates a tension identified across multiple theoretical traditions.

Consider Lin and Ma's multirelation approach to general systems \cite{Lin1987}, wherein a system $S = (M, R)$ consists of an object set $M$ and relation set $R$. This formulation appears isomorphic to the extensional relational structure, yet Lin emphasizes that ``the concept of systems is a generalization of that of structures.'' The crucial distinction emerges when we recognize that Lin's ``system'' already incorporates the temporal and modal dimensions that Guarino's extensional structure deliberately brackets.

Mesarovic and Takahara's definition—``a system $S$ is a relation on nonempty sets'' \cite{Mesarovic1975}—further exemplifies this minimalist extensional approach:
$$S \subseteq \prod\{V_i : i \in I\}$$
Here the Cartesian product structure makes explicit what remains implicit in Guarino's formulation: the arity and type structure of relations themselves become first-class ontological citizens.

\subsection{The Problem of Summative versus Constitutive Characteristics}

Von Bertalanffy's distinction between summative and constitutive properties illuminates a critical limitation of purely extensional structures. He writes: ``Summative characteristics of an element are those which are the same within and outside the complex... Constitutive characteristics are those which are dependent on the specific relations within the complex'' \cite{Bertalanffy1968}.

An extensional structure $(D, \mathbf{R})$ captures relations as sets of tuples, but it cannot inherently distinguish between:
\begin{enumerate}
	\item Relations that are \emph{intrinsic} to the system's organization (constitutive)
	\item Relations that merely happen to hold (summative aggregations)
\end{enumerate}

\begin{example}[Molecular Isomerism as Relational Complexity]
	Consider two chemical compounds with identical molecular formulas but different structural arrangements (isomers). In the extensional framework:
	\begin{align*}
		D &= \{\text{atom}_1, \text{atom}_2, \ldots, \text{atom}_n\} \\
		\mathbf{R}_{\text{compound1}} &= \{\text{bonds-to}, \text{shares-electrons}, \ldots\} \\
		\mathbf{R}_{\text{compound2}} &= \{\text{bonds-to}', \text{shares-electrons}', \ldots\}
	\end{align*}
	
	The two structures $(D, \mathbf{R}_{\text{compound1}})$ and $(D, \mathbf{R}_{\text{compound2}})$ share the same universe $D$ but differ in their relational extensions. The emergent chemical properties—different boiling points, reactivities—exemplify constitutive characteristics that arise from these specific relational configurations. Yet the extensional structure alone provides no principle for distinguishing these constitutive relations from mere accidental co-occurrences.
\end{example}